\section{Ubicación de la clase: entender Mixtral exige llevar tres cuentas a la vez}

Mixtral suele resumirse en una sola frase: «ocho expertos de 7B, y cada token activa solo dos». La frase basta para atraer la atención, pero no explica por qué el modelo es rápido, por qué ocupa tanta memoria de GPU ni por qué necesita una comunicación compleja, y tampoco aclara qué aprende en realidad cada «experto». El punto de entrada correcto de esta clase es llevar tres cuentas a la vez: el \textbf{capacidad total} registra cuántos parámetros pueden almacenar todos los expertos; el \textbf{cómputo activo} registra cuántos parámetros recorre realmente un token; y el \textbf{costo del sistema} registra la residencia de los pesos, el token dispatch, la comunicación entre tarjetas y el desequilibrio de carga. Ninguna de las tres sustituye a las otras.

\subsection{Primero se corrige la fuente; después se traza la ruta de lectura}

Esta sección aclara primero los límites del material, porque el directorio anterior tenía una contaminación grave de fuentes: el archivo de diapositivas anterior era en realidad el deck de la Lecture 28 de Nathan Lambert sobre alineación, no la clase de Albert Jiang sobre Mixtral. El sitio oficial de CS25 no publicó un deck propio para esta clase, así que aquí se reconstruyen los slide states segundo a segundo a partir del video oficial en 1920x1080, y los 34 minutos de Q\&A se incorporan al texto como teacher voice. Todas las figuras que siguen provienen del video oficial, no del PDF equivocado.

\lecturefigure{slide-01-content.jpg}{Ruta de la clase: dense Transformer, Sparse Mixture of Experts y routing interpretation}{Video oficial de Stanford Online 00:01:59}

\begin{importantbox}{La pregunta central de esta clase}
Sparse Mixture of Experts (SMoE, mezcla dispersa de expertos) no es «un dense model más grande». Amplía el único MLP de cada capa a varios expert MLP, y un router elige un número reducido de expertos para cada token. Así se puede aumentar la capacidad total sin aumentar en la misma proporción los FLOPs por token; sin embargo, todos los expertos siguen teniendo que almacenarse, los tokens siguen teniendo que moverse entre dispositivos y el router todavía puede concentrar la carga en unos pocos expertos.
\end{importantbox}

\begin{knowledgebox}{Términos clave: qué responde cada una de las tres cuentas}
\begin{tabularx}{\textwidth}{L{0.18\textwidth} L{0.28\textwidth} X}
\toprule
Cuenta & Pregunta central & Variables típicas \\
\midrule
Capacidad total & ¿Cuántos pesos almacena el modelo en total? & total parameters, todos los experts, shared attention \\
Cómputo activo & ¿Qué pesos recorre realmente un token? & top-$k$, active parameters, per-token FLOPs \\
Costo del sistema & ¿Cómo se distribuyen los pesos y los tokens en el hardware? & HBM (High Bandwidth Memory, memoria DRAM de alto ancho de banda), expert parallelism, all-to-all, lote, load balance \\
\bottomrule
\end{tabularx}
\end{knowledgebox}

\teachervoice{00:01:04--00:01:54. Albert divide explícitamente la clase en dos partes, architecture e interpretation, e intercala una y otra vez open research questions. Espera que la comunidad abierta siga investigando, en lugar de tratar a Mixtral como una receta ya explicada por completo.}

\subsection{Niveles de evidencia: qué es medición y qué es solo hipótesis de trabajo}

Esta clase reúne configuraciones del modelo, benchmark plots, routing histograms, ablations de la comunidad y juicios basados en la experiencia expresados en el Q\&A. Al leer hay que distinguirlos: la tabla de parámetros del modelo es un hecho estructural; los benchmarks de 2024 son evidencia comparativa de ese momento; «el MLP almacena conocimiento y la attention razona» es la conventional wisdom que el speaker enuncia de forma explícita; que cierto expert se elija con más frecuencia en datos de matemáticas es solo una correlación; y «un expert count mayor facilita la especialización» es un juicio prospectivo. Más adelante se marca el límite en cada caso.

\begin{warningbox}{El expert del nombre no equivale a un especialista profesional humano}
Un expert es solo una subred de parámetros que el router puede seleccionar. Puede ser sensible a la forma de los tokens, a la sintaxis local, a la posición, a la frecuencia o a combinaciones lineales de varias características latentes, y no corresponde de manera natural a un «experto en código» o un «experto en medicina». Traducir directamente la selección del router a semántica humana es el cuarto myth que esta clase busca corregir.
\end{warningbox}

\subsection{Resumen de la sección}

El hilo didáctico de esta clase no es la promoción del modelo, sino reunir sparse capacity, active compute, resident memory, communication e interpretabilidad en un mismo sistema. La corrección de la fuente implica además que toda la slide coverage y las explicaciones orales del docente deben reconstruirse; las notas anteriores no pueden reutilizarse.
